% ============================================================
%  يُترجَم بـ XeLaTeX  (في Overleaf: Menu → Compiler → XeLaTeX)
%  الخط: Amiri (متوفر في Overleaf، أو حمّله من github.com/aliftype/amiri)
% ============================================================
\documentclass[12pt,a4paper]{article}
\usepackage[a4paper,margin=1.7cm,top=2.3cm,bottom=2.9cm,headheight=15pt,footskip=1.5cm]{geometry}
\usepackage{amsmath,amssymb}
\usepackage{icomma}
\usepackage[english,bidi=default,layout=lists]{babel}
\babelprovide[import,main]{arabic}
\babelfont{rm}[Scale=1.08]{Amiri}
\babelfont{sf}[Scale=1.08]{Amiri}
\usepackage[most]{tcolorbox}
\usepackage{tikz}
\usetikzlibrary{arrows.meta,calc,shadings}
\usepackage{array}
\usepackage{colortbl}
\usepackage{fancyhdr}
\usepackage{titlesec}
\usepackage{enumitem}
\usepackage{collcell}
\usepackage{needspace}
\usepackage{etoolbox}
\newcolumntype{C}{>{\collectcell\ar}c<{\endcollectcell}}
\pretocmd{\section}{\Needspace{15\baselineskip}}{}{}
\pretocmd{\subsection}{\Needspace{11\baselineskip}}{}{}

\definecolor{cDef}{HTML}{1F4E9C}
\definecolor{cThm}{HTML}{B3261E}
\definecolor{cRem}{HTML}{2E7D32}
\definecolor{cEx}{HTML}{C76B00}
\definecolor{cMeth}{HTML}{6A1B9A}
\definecolor{cSol}{HTML}{00695C}
\definecolor{cSec}{HTML}{1F3A68}

\titleformat{\section}{\Large\bfseries\color{cSec}}{\thesection.}{0.6em}{}[{\color{cSec!60}\titlerule[1pt]}]
\titleformat{\subsection}{\large\bfseries\color{cThm!85!black}}{\thesubsection.}{0.5em}{}
\titlespacing*{\section}{0pt}{14pt}{8pt}
\titlespacing*{\subsection}{0pt}{10pt}{5pt}
\setlength{\parindent}{0pt}
\setlength{\parskip}{4pt}
\setlist{itemsep=2pt,topsep=3pt}

\newtcolorbox{arbox}[3][]{enhanced,
  colback=#2!5,colframe=#2!80!black,arc=2mm,boxrule=0.8pt,
  left=6pt,right=6pt,top=8pt,bottom=5pt,
  fonttitle=\bfseries,coltitle=white,
  title={\foreignlanguage{arabic}{#3}},
  attach boxed title to top right={xshift=-5mm,yshift=-3mm},
  boxed title style={colback=#2!80!black,arc=1.5mm,boxrule=0pt},
  before upper={\selectlanguage{arabic}},
  before skip=8pt,after skip=8pt,breakable,#1}
\newenvironment{BOX}[2]{\begin{otherlanguage}{english}\begin{arbox}{#1}{#2}}{\end{arbox}\end{otherlanguage}}
\newenvironment{tarif}[1][تعريف]{\begin{BOX}{cDef}{#1}}{\end{BOX}}
\newenvironment{khassia}[1][خاصية]{\begin{BOX}{cThm}{#1}}{\end{BOX}}
\newenvironment{molahada}[1][ملاحظة]{\begin{BOX}{cRem}{#1}}{\end{BOX}}
\newenvironment{mithal}[1][مثال]{\begin{BOX}{cEx}{#1}}{\end{BOX}}
\newenvironment{tariqa}[1][طريقة]{\begin{BOX}{cMeth}{#1}}{\end{BOX}}

\newcounter{exo}
\newenvironment{tamrin}[1]{\stepcounter{exo}\begin{BOX}{cSec}{التمرين \theexo: #1}}{\end{BOX}}
\newenvironment{hal}{\Needspace{12\baselineskip}\begin{otherlanguage}{english}\begin{arbox}[breakable]{cSol}{حل التمرين \theexo}}{\end{arbox}\end{otherlanguage}}

\newenvironment{LTR}{\begin{otherlanguage}{english}}{\end{otherlanguage}}
\newcommand{\ar}[1]{\foreignlanguage{arabic}{#1}}

\renewcommand{\dots}{\mathinner{\ldotp\ldotp\ldotp}}
\renewcommand{\ldots}{\mathinner{\ldotp\ldotp\ldotp}}
\renewcommand{\cdots}{\mathinner{\cdotp\cdotp\cdotp}}
\newcommand{\N}{\mathbb{N}}
\newcommand{\Z}{\mathbb{Z}}
\newcommand{\D}{\mathbb{D}}
\newcommand{\Q}{\mathbb{Q}}
\newcommand{\R}{\mathbb{R}}
\newcommand{\pgcd}{\operatorname{PGCD}}
\newcommand{\ppcm}{\operatorname{PPCM}}

% \NLseg{a}{b}{[}{]} ; \NLright{a}{[} ; \NLleft{b}{]}  (interval on a number line)
\newcommand{\NLseg}[4]{\begin{LTR}\begin{tikzpicture}[baseline=-0.6ex]
 \draw[-{Stealth[length=2mm]},gray!70](0,0)--(3.4,0);
 \draw[line width=2pt,cThm!80!black](0.7,0)--(2.7,0);
 \node[cThm!80!black,font=\LARGE] at (0.7,0){$#3$};
 \node[cThm!80!black,font=\LARGE] at (2.7,0){$#4$};
 \node[below=3pt,font=\footnotesize] at (0.7,0){$#1$};
 \node[below=3pt,font=\footnotesize] at (2.7,0){$#2$};
\end{tikzpicture}\end{LTR}}
\newcommand{\NLright}[2]{\begin{LTR}\begin{tikzpicture}[baseline=-0.6ex]
 \draw[-{Stealth[length=2mm]},gray!70](0,0)--(3.4,0);
 \draw[line width=2pt,cThm!80!black](1.3,0)--(3.3,0);
 \node[cThm!80!black,font=\LARGE] at (1.3,0){$#2$};
 \node[below=3pt,font=\footnotesize] at (1.3,0){$#1$};
\end{tikzpicture}\end{LTR}}
\newcommand{\NLleft}[2]{\begin{LTR}\begin{tikzpicture}[baseline=-0.6ex]
 \draw[-{Stealth[length=2mm]},gray!70](0,0)--(3.4,0);
 \draw[line width=2pt,cThm!80!black](0.05,0)--(2.1,0);
 \node[cThm!80!black,font=\LARGE] at (2.1,0){$#2$};
 \node[below=3pt,font=\footnotesize] at (2.1,0){$#1$};
\end{tikzpicture}\end{LTR}}

% ---------- social-media icons (pure TikZ) ----------
\definecolor{cFB}{HTML}{1877F2}
\definecolor{cTG}{HTML}{229ED9}
\definecolor{cIG1}{HTML}{FEDA75}
\definecolor{cIG2}{HTML}{D62976}
\definecolor{cIG3}{HTML}{4F5BD5}
\newcommand{\fbicon}{\begin{LTR}\begin{tikzpicture}[baseline=-0.75ex,x=1em,y=1em]
 \fill[cFB] (0,0) circle (0.55);
 \node[white,inner sep=0,font=\usefont{TU}{lmss}{bx}{n}\small] at (0.04,-0.02) {f};
\end{tikzpicture}\end{LTR}}
\newcommand{\igicon}{\begin{LTR}\begin{tikzpicture}[baseline=-0.75ex,x=1em,y=1em]
 \shade[lower left=cIG1,lower right=cIG2,upper left=cIG2,upper right=cIG3,rounded corners=0.28em] (-0.55,-0.55) rectangle (0.55,0.55);
 \draw[white,line width=0.09em,rounded corners=0.18em] (-0.36,-0.36) rectangle (0.36,0.36);
 \draw[white,line width=0.09em] (0,0) circle (0.17);
 \fill[white] (0.22,0.22) circle (0.05);
\end{tikzpicture}\end{LTR}}
\newcommand{\tgicon}{\begin{LTR}\begin{tikzpicture}[baseline=-0.75ex,x=1em,y=1em]
 \fill[cTG] (0,0) circle (0.55);
 \fill[white] (-0.33,0.0) -- (0.32,0.27) -- (0.2,-0.3) -- (0.0,-0.15) -- (-0.08,-0.28) -- (-0.08,-0.08) -- cycle;
 \fill[cTG!70!black] (-0.08,-0.08) -- (0.25,0.2) -- (0.0,-0.15) -- (-0.08,-0.28) -- cycle;
\end{tikzpicture}\end{LTR}}
% page number: black box, white two-digit number, at the right of the social row
\newcommand{\PageBox}{\begin{LTR}\tikz[baseline=(p.base)]\node[fill=black,text=white,inner sep=2.5pt,minimum width=1.6em,minimum height=1.3em,font=\usefont{TU}{lmss}{bx}{n}\small](p){\ifnum\value{page}<10 0\fi\arabic{page}};\end{LTR}}
% visual left-to-right order: Telegram | Instagram | Facebook | page box
\newcommand{\SocialRow}{\begin{otherlanguage}{english}%
 \ar{\textbf{الأستاذ سوايسية هشام للرياضيات}}\ \tgicon\hfill
 \textbf{prof\_hicham\_math}\ \igicon\hfill
 \ar{\textbf{الأستاذ سوايسية هشام للرياضيات}}\ \fbicon\hspace{2.2em}\PageBox\end{otherlanguage}}

\pagestyle{fancy}
\fancyhf{}
\renewcommand{\headrulewidth}{0.6pt}
\renewcommand{\footrulewidth}{0.4pt}
% \SetHeader{right head: document title}{left head: level -- stream}
\newcommand{\SetHeader}[2]{%
 \fancyhead[C]{\parbox{\textwidth}{\ar{#1\hfill #2}}}%
 \fancyfoot[C]{\parbox{\textwidth}{\small\SocialRow}}}
\usepackage{tkz-tab}
\usepackage{multicol}
% variation / sign table (LTR)
\newenvironment{TAB}{\par\begin{center}\begin{LTR}\begin{tikzpicture}[scale=0.95,every node/.style={font=\small}]}{\end{tikzpicture}\end{LTR}\end{center}}
% grid + axes, clipped: \begin{GR}{xmin}{xmax}{ymin}{ymax}{unit cm}
\newenvironment{GR}[5]{\par\begin{center}\begin{LTR}\begin{tikzpicture}[x=#5cm,y=#5cm,font=\footnotesize]
 \draw[gray!25,very thin] (#1,#3) grid (#2,#4);
 \draw[-{Stealth[length=2mm]}] (#1,0)--(#2+0.3,0) node[below left]{$x$};
 \draw[-{Stealth[length=2mm]}] (0,#3)--(0,#4+0.3) node[below left]{$y$};
 \node[below left] at (0,0){$O$};
 \begin{scope}\clip (#1,#3) rectangle (#2,#4);}
{\end{scope}\end{tikzpicture}\end{LTR}\end{center}}
\newcommand{\Ticks}[2]{\foreach \t in {#1} \node[below,font=\scriptsize] at (\t,0){$\t$}; \foreach \t in {#2} \node[left,font=\scriptsize] at (0,\t){$\t$};}
\newcommand{\bx}[1]{\ensuremath{\colorbox{cDef!12}{$\displaystyle #1$}}}
\newcommand{\TwoCol}[4]{\par\begin{LTR}\noindent\begin{minipage}[t]{#1\textwidth}\selectlanguage{arabic}#2\end{minipage}\hfill\begin{minipage}[t]{#3\textwidth}\selectlanguage{arabic}#4\end{minipage}\end{LTR}\par}
\newcommand{\RL}[4]{\TwoCol{#3}{#4}{#1}{#2}}  % \RL{w1}{right content}{w2}{left content}
\newcommand{\ok}{$\textcolor{cRem}{\checkmark}$}
\newcommand{\limn}{\displaystyle\lim_{n\to+\infty}}
% sub-questions, each on its own line
\newenvironment{sq}{\begin{itemize}[leftmargin=2.4em,itemsep=3pt]}{\end{itemize}}
\newcommand{\qa}{\item[(أ)]}\newcommand{\qb}{\item[(ب)]}\newcommand{\qc}{\item[(جـ)]}\newcommand{\qd}{\item[(د)]}\newcommand{\qe}{\item[(هـ)]}
\newcommand{\pa}{\par\smallskip\textbf{(أ)} }\newcommand{\pb}{\par\smallskip\textbf{(ب)} }\newcommand{\pc}{\par\smallskip\textbf{(جـ)} }\newcommand{\pd}{\par\smallskip\textbf{(د)} }\newcommand{\pe}{\par\smallskip\textbf{(هـ)} }
\newcommand{\QQ}[1]{\par\smallskip\textbf{#1.} }
\newenvironment{tmr}{\stepcounter{exo}\begin{BOX}{cSec}{تمرين محلول \theexo}}{\end{BOX}}
\newcommand{\ABC}[3]{\\[2pt]\begin{tabular}{p{0.28\textwidth}p{0.28\textwidth}p{0.28\textwidth}}(جـ) \ #3 & (ب) \ #2 & (أ) \ #1\end{tabular}}
\pretocmd{\section}{\Needspace{24\baselineskip}}{}{}
\newtcbox{\bxtbox}{on line,colback=cDef!12,colframe=cDef!12,boxrule=0pt,arc=1pt,boxsep=0pt,left=2pt,right=2pt,top=1pt,bottom=1pt}
\newcommand{\bxt}[1]{\foreignlanguage{english}{\bxtbox{\ar{#1}}}}
\SetHeader{\small\bfseries\color{cSec}تمارين البكالوريا: الدالة الأسية}{\small\color{cSec}السنة الثالثة ثانوي -- الشعب العلمية}
\newenvironment{FIG}{\par\begin{center}\begin{LTR}}{\end{LTR}\end{center}}
\begin{document}
\begin{LTR}
\begin{tcolorbox}[enhanced,colback=cSec,colframe=cSec,arc=3mm,halign=center,boxrule=0pt,top=12pt,bottom=12pt,before upper={\selectlanguage{arabic}}]
{\color{white}\fontsize{26}{32}\selectfont\bfseries تمارين البكالوريا في الدالة الأسية}\\[8pt]
{\color{white}\large عيّنة مجانية: موضوعا 2026 و 2025 بالحلول المفصلة}\\[5pt]
{\color{white}\large علوم تجريبية -- تقني رياضي -- رياضيات}
\end{tcolorbox}
\end{LTR}
\section{تمارين البكالوريا -- شعبة علوم تجريبية}
% ===================== بكالوريا 2026 -- الموضوع الثاني =====================
\begin{tamrin}{بكالوريا 2026 -- الموضوع الثاني -- علوم تجريبية}
$f$ الدالة المعرفة على $\R^{*}$ بـ: $f(x)=1+\dfrac{5-4e^{x}}{e^{2x}-1}$، و $(C_f)$ تمثيلها البياني في المستوي المنسوب إلى المعلم المتعامد والمتجانس $(O;\vec i,\vec j)$.
\QQ{1} أحسب النهايات: $\lim\limits_{x\to-\infty}f(x)$، $\lim\limits_{x\to0^{+}}f(x)$، $\lim\limits_{x\to0^{-}}f(x)$ و $\lim\limits_{x\to+\infty}f(x)$، ثم فسّر النتائج هندسيا.
\QQ{2}
\begin{sq}
\qa بيّن أنه من أجل كل $x$ من $\R^{*}$: $f'(x)=\dfrac{2e^{x}\left(e^{x}-2\right)\left(2e^{x}-1\right)}{\left(e^{2x}-1\right)^{2}}$.
\qb استنتج اتجاه تغيّر الدالة $f$ ثم شكّل جدول تغيّراتها.
\end{sq}
\QQ{3} بيّن أنه من أجل كل $x$ من $\R^{*}$: $f(-x)+f(x)=-3$، وفسّر النتيجة هندسيا، ثم أرسم المنحنى $(C_f)$.
\QQ{4} تحقّق أنه من أجل كل $x$ من $\R^{*}$: $f(x)=-4+\dfrac{1}{2}\left(\dfrac{9e^{x}}{1+e^{x}}-\dfrac{e^{x}}{1-e^{x}}\right)$.
\QQ{5} $g$ الدالة المعرفة على $]1;+\infty[$ بـ: $g(x)=f(\ln x)$.
تحقّق أنه من أجل كل $x>1$: $g(x)=\dfrac{(x-2)^{2}}{x^{2}-1}$، وحدّد اتجاه تغيّر الدالة $g$ ثم شكّل جدول تغيّراتها.
\end{tamrin}


\begin{hal}
\QQ{1} \textbf{حساب النهايات:}
\begin{itemize}
\item عند $-\infty$: $\lim\limits_{x\to-\infty}e^{x}=0$ و $\lim\limits_{x\to-\infty}e^{2x}=0$، إذن $\lim\limits_{x\to-\infty}\dfrac{5-4e^{x}}{e^{2x}-1}=\dfrac{5}{-1}=-5$، ومنه:
\[\bx{\lim_{x\to-\infty}f(x)=1-5=-4}\]
\item عند $0$: $\lim\limits_{x\to0}(5-4e^{x})=5-4=1$، و $e^{2x}-1>0$ إذا كان $x>0$، و $e^{2x}-1<0$ إذا كان $x<0$، إذن $\lim\limits_{x\to0^{+}}(e^{2x}-1)=0^{+}$ و $\lim\limits_{x\to0^{-}}(e^{2x}-1)=0^{-}$، ومنه:
\[\bx{\lim_{x\to0^{+}}f(x)=+\infty}\qquad \bx{\lim_{x\to0^{-}}f(x)=-\infty}\]
\item عند $+\infty$: نقسم البسط والمقام على $e^{2x}$: $\dfrac{5-4e^{x}}{e^{2x}-1}=\dfrac{5e^{-2x}-4e^{-x}}{1-e^{-2x}}$، وبما أن $\lim\limits_{x\to+\infty}e^{-x}=\lim\limits_{x\to+\infty}e^{-2x}=0$ فإن هذا الكسر يؤول إلى $\dfrac{0}{1}=0$، ومنه:
\[\bx{\lim_{x\to+\infty}f(x)=1}\]
\end{itemize}
\textbf{التفسير الهندسي:}
\begin{itemize}
\item المستقيم ذو المعادلة $y=-4$ مقارب أفقي لـ $(C_f)$ بجوار $-\infty$.
\item المستقيم ذو المعادلة $x=0$ (محور التراتيب) مقارب عمودي لـ $(C_f)$.
\item المستقيم ذو المعادلة $y=1$ مقارب أفقي لـ $(C_f)$ بجوار $+\infty$.
\end{itemize}

\QQ{2}
\pa الدالة $f$ قابلة للاشتقاق على $\R^{*}$، ومن أجل كل $x\neq0$ (مشتقة حاصل القسمة):
\begin{align*}
f'(x)&=\frac{-4e^{x}\left(e^{2x}-1\right)-\left(5-4e^{x}\right)2e^{2x}}{\left(e^{2x}-1\right)^{2}}
=\frac{-4e^{3x}+4e^{x}-10e^{2x}+8e^{3x}}{\left(e^{2x}-1\right)^{2}}\\
&=\frac{4e^{3x}-10e^{2x}+4e^{x}}{\left(e^{2x}-1\right)^{2}}=\frac{2e^{x}\left(2e^{2x}-5e^{x}+2\right)}{\left(e^{2x}-1\right)^{2}}
\end{align*}
ولدينا $\left(e^{x}-2\right)\left(2e^{x}-1\right)=2e^{2x}-e^{x}-4e^{x}+2=2e^{2x}-5e^{x}+2$، إذن:
\[\bx{f'(x)=\frac{2e^{x}\left(e^{x}-2\right)\left(2e^{x}-1\right)}{\left(e^{2x}-1\right)^{2}}}\]

\pb بما أن $2e^{x}>0$ و $\left(e^{2x}-1\right)^{2}>0$ على $\R^{*}$، فإن إشارة $f'(x)$ هي إشارة الجداء $\left(e^{x}-2\right)\left(2e^{x}-1\right)$:
\begin{itemize}
\item $e^{x}-2>0 \iff e^{x}>2 \iff x>\ln2$.
\item $2e^{x}-1>0 \iff e^{x}>\dfrac12 \iff x>-\ln2$.
\end{itemize}
القيم الحدية: $f(\ln2)=1+\dfrac{5-8}{4-1}=1-1=0$ و $f(-\ln2)=1+\dfrac{5-2}{\frac14-1}=1+\dfrac{3}{-\frac34}=1-4=-3$.

إذن: \bxt{$f$ متزايدة تماما على كل من $]{-}\infty;-\ln2]$ و $[\ln2;+\infty[$}

و \bxt{$f$ متناقصة تماما على كل من $[-\ln2;0[$ و $]0;\ln2]$}.
\begin{TAB}
\tkzTabInit[lgt=2,espcl=2.4]{$x$/1,$e^{x}-2$/1,$2e^{x}-1$/1,$f'(x)$/1,$f(x)$/2.2}{$-\infty$,$-\ln2$,$0$,$\ln2$,$+\infty$}
\tkzTabLine{,-,t,-,d,-,z,+,}
\tkzTabLine{,-,z,+,d,+,t,+,}
\tkzTabLine{,+,z,-,d,-,z,+,}
\tkzTabVar{-/$-4$,+/$-3$,-D+/$-\infty$/$+\infty$,-/$0$,+/$1$}
\end{TAB}

\QQ{3} من أجل كل $x\in\R^{*}$ لدينا $-x\in\R^{*}$، و:
\[f(-x)=1+\frac{5-4e^{-x}}{e^{-2x}-1}=1+\frac{e^{2x}\left(5-4e^{-x}\right)}{e^{2x}\left(e^{-2x}-1\right)}=1+\frac{5e^{2x}-4e^{x}}{1-e^{2x}}\]
ومنه:
\begin{align*}
f(-x)+f(x)&=2+\frac{5-4e^{x}}{e^{2x}-1}-\frac{5e^{2x}-4e^{x}}{e^{2x}-1}
=2+\frac{5-5e^{2x}}{e^{2x}-1}=2-5
\end{align*}
إذن \bx{f(-x)+f(x)=-3}.

\textbf{التفسير الهندسي:} المساواة من الشكل $f(2a-x)+f(x)=2b$ مع $a=0$ و $b=-\dfrac32$، إذن النقطة \bx{\Omega\left(0;-\tfrac32\right)} مركز تناظر للمنحنى $(C_f)$.

\textbf{الرسم:} نلاحظ أن $f(x)=\dfrac{e^{2x}-1+5-4e^{x}}{e^{2x}-1}=\dfrac{\left(e^{x}-2\right)^{2}}{e^{2x}-1}$، إذن $(C_f)$ يقطع محور الفواصل في النقطة $(\ln2;0)$ فقط (ويمسّه فيها لأن $f'(\ln2)=0$)، ولا يقطع محور التراتيب. نعتمد: $\ln2\approx0{,}69$، والذروة $(-\ln2;-3)$.
\begin{GR}{-4.5}{4.5}{-7}{3}{0.85}
\Ticks{-4,-3,-2,-1,1,2,3,4}{-6,-5,-4,-3,-2,-1,1,2}
\draw[cThm,dashed,thick](-4.5,1)--(4.5,1) node[pos=0.9,above]{$y=1$};
\draw[cThm,dashed,thick](-4.5,-4)--(4.5,-4) node[pos=0.12,below]{$y=-4$};
\draw[cDef,very thick,domain=-4.5:-0.075,samples=200] plot (\x,{1+(5-4*exp(\x))/(exp(2*\x)-1)});
\draw[cDef,very thick,domain=0.075:4.5,samples=200] plot (\x,{1+(5-4*exp(\x))/(exp(2*\x)-1)});
\fill[cEx](0,-1.5) circle(2pt) node[left]{$\Omega$};
\fill(0.693,0) circle(1.5pt) node[above]{$\ln2$};
\fill(-0.693,-3) circle(1.5pt);
\draw[densely dotted](-0.693,0)--(-0.693,-3);
\node[cDef] at (1.6,2.4){$(C_f)$};
\end{GR}

\QQ{4} نوحّد المقامات، مع $(1+e^{x})(1-e^{x})=1-e^{2x}$:
\[\frac{9e^{x}}{1+e^{x}}-\frac{e^{x}}{1-e^{x}}=\frac{9e^{x}(1-e^{x})-e^{x}(1+e^{x})}{1-e^{2x}}=\frac{8e^{x}-10e^{2x}}{1-e^{2x}}\]
ومنه:
\begin{align*}
-4+\frac12\left(\frac{9e^{x}}{1+e^{x}}-\frac{e^{x}}{1-e^{x}}\right)&=-4+\frac{4e^{x}-5e^{2x}}{1-e^{2x}}
=\frac{-4+4e^{2x}+4e^{x}-5e^{2x}}{1-e^{2x}}\\&=\frac{-e^{2x}+4e^{x}-4}{1-e^{2x}}=\frac{e^{2x}-4e^{x}+4}{e^{2x}-1}
\end{align*}
ومن جهة أخرى $f(x)=\dfrac{e^{2x}-1+5-4e^{x}}{e^{2x}-1}=\dfrac{e^{2x}-4e^{x}+4}{e^{2x}-1}$. إذن:
\[\bx{f(x)=-4+\frac{1}{2}\left(\frac{9e^{x}}{1+e^{x}}-\frac{e^{x}}{1-e^{x}}\right)}\]

\QQ{5} من أجل $x>1$ لدينا $\ln x>0$ أي $\ln x\in\R^{*}$، إذن $g(x)=f(\ln x)$ معرفة، و $e^{\ln x}=x$، $e^{2\ln x}=x^{2}$:
\[g(x)=1+\frac{5-4x}{x^{2}-1}=\frac{x^{2}-1+5-4x}{x^{2}-1}=\frac{x^{2}-4x+4}{x^{2}-1}\]
إذن \bx{g(x)=\dfrac{(x-2)^{2}}{x^{2}-1}}.

\textbf{اتجاه التغيّر:} $g=f\circ\ln$، و $g'(x)=\dfrac{1}{x}\,f'(\ln x)$، و $\dfrac1x>0$، إذن إشارة $g'(x)$ هي إشارة $f'(\ln x)$:
\begin{itemize}
\item إذا كان $1<x<2$ فإن $0<\ln x<\ln2$، و $f'<0$ على $]0;\ln2[$، إذن $g'(x)<0$.
\item إذا كان $x>2$ فإن $\ln x>\ln2$، و $f'>0$ على $]\ln2;+\infty[$، إذن $g'(x)>0$.
\end{itemize}
إذن \bxt{$g$ متناقصة تماما على $]1;2]$ ومتزايدة تماما على $[2;+\infty[$}.

\textbf{النهايات:} $\lim\limits_{x\to1^{+}}\ln x=0^{+}$ و $\lim\limits_{t\to0^{+}}f(t)=+\infty$، إذن $\lim\limits_{x\to1^{+}}g(x)=+\infty$. و $\lim\limits_{x\to+\infty}\ln x=+\infty$ و $\lim\limits_{t\to+\infty}f(t)=1$، إذن $\lim\limits_{x\to+\infty}g(x)=1$. و $g(2)=0$.
\begin{TAB}
\tkzTabInit[lgt=2,espcl=3]{$x$/1,$g'(x)$/1,$g(x)$/2}{$1$,$2$,$+\infty$}
\tkzTabLine{d,-,z,+,}
\tkzTabVar{D+/$+\infty$,-/$0$,+/$1$}
\end{TAB}
\end{hal}

% ===================== بكالوريا 2025 -- الموضوع الأول =====================
\begin{tamrin}{بكالوريا 2025 -- الموضوع الأول -- علوم تجريبية}
$f$ الدالة المعرفة على $\R$ بـ: $f(x)=e^{2x}-e^{x}-x-2$، و $(C_f)$ تمثيلها البياني في المستوي المنسوب إلى المعلم المتعامد والمتجانس $(O;\vec i,\vec j)$ (وحدة الطول $2cm$).
\QQ{1}
\begin{sq}
\qa أحسب $\lim\limits_{x\to-\infty}f(x)$ ثم بيّن أن $\lim\limits_{x\to+\infty}f(x)=+\infty$.
\qb بيّن أن المستقيم $(\Delta)$ ذا المعادلة $y=-x-2$ مقارب مائل لـ $(C_f)$ عند $-\infty$.
\qc أدرس الوضع النسبي للمنحنى $(C_f)$ والمستقيم $(\Delta)$.
\end{sq}
\QQ{2}
\begin{sq}
\qa بيّن أنه من أجل كل عدد حقيقي $x$: $f'(x)=\left(e^{x}-1\right)\left(2e^{x}+1\right)$.
\qb استنتج اتجاه تغيّر الدالة $f$ ثم شكّل جدول تغيّراتها.
\end{sq}
\QQ{3} بيّن أن المنحنى $(C_f)$ يقبل مماسا $(T)$ موازيا للمستقيم $(\Delta)$، يُطلب تعيين معادلة له.
\QQ{4} بيّن أن المعادلة $f(x)=0$ تقبل حلّين فقط $\alpha$ و $\beta$، ثم تحقّق أن: $-2{,}2<\alpha<-2{,}1$ و $0{,}8<\beta<0{,}9$.
\QQ{5} بيّن أن المنحنى $(C_f)$ يقبل نقطة انعطاف يُطلب تعيين إحداثيّيها.
\QQ{6}
\begin{sq}
\qa أرسم كلا من $(\Delta)$، $(T)$ و $(C_f)$.
\qb ناقش بيانيا، وحسب قيم الوسيط الحقيقي $m$، عدد حلول المعادلة: $e^{2x}-e^{x}-m-2=0$.
\end{sq}
\end{tamrin}

\begin{hal}
\QQ{1}
\pa عند $-\infty$: $\lim\limits_{x\to-\infty}e^{2x}=\lim\limits_{x\to-\infty}e^{x}=0$ و $\lim\limits_{x\to-\infty}(-x-2)=+\infty$، إذن:
\[\bx{\lim_{x\to-\infty}f(x)=+\infty}\]
عند $+\infty$ (حالة عدم تعيين $+\infty-\infty$): نخرج $e^{2x}$ عاملا مشتركا:
\[f(x)=e^{2x}\left(1-e^{-x}-xe^{-2x}-2e^{-2x}\right)\]
لدينا $\lim\limits_{x\to+\infty}e^{-x}=\lim\limits_{x\to+\infty}e^{-2x}=0$ و $\lim\limits_{x\to+\infty}xe^{-2x}=\lim\limits_{x\to+\infty}\left(xe^{-x}\right)e^{-x}=0$ (التزايد المقارن: $\lim\limits_{x\to+\infty}xe^{-x}=0$)، إذن القوس يؤول إلى $1$، و $\lim\limits_{x\to+\infty}e^{2x}=+\infty$، ومنه \bx{\lim\limits_{x\to+\infty}f(x)=+\infty}.

\pb $f(x)-(-x-2)=e^{2x}-e^{x}$، و $\lim\limits_{x\to-\infty}\left(e^{2x}-e^{x}\right)=0$، إذن \bxt{المستقيم $(\Delta):y=-x-2$ مقارب مائل لـ $(C_f)$ عند $-\infty$}.

\pc $f(x)-(-x-2)=e^{x}\left(e^{x}-1\right)$، و $e^{x}>0$، فالإشارة هي إشارة $e^{x}-1$، و $e^{x}-1>0\iff x>0$:
\begin{FIG}
\begin{tikzpicture}[font=\normalsize]
\fill[gray!12] (0,0) rectangle (2.6,-0.9);
\fill[green!18] (2.6,0) rectangle (15.5,-0.9);
\fill[orange!45] (0,-0.9) rectangle (2.6,-4.6);
\fill[cyan!22] (2.6,-0.9) rectangle (15.5,-4.6);
\draw[thick] (0,0) rectangle (15.5,-4.6); \draw (0,-0.9)--(15.5,-0.9) (0,-1.8)--(15.5,-1.8) (2.6,0)--(2.6,-4.6);
\node at (1.3,-0.45){$x$};
\node at (1.3,-1.35){$f(x)-(-x-2)$};
\node[align=center,font=\bfseries] at (1.3,-3.2){\ar{الوضع}\\ \ar{النسبي}};
\node at (3.3,-0.45){$-\infty$};
\node at (9.05,-0.45){$0$};
\node at (14.8,-0.45){$+\infty$};
\node at (5.425,-1.35){$-$};
\node[align=center,font=\small] at (5.425,-3.2){$(C_f)$\\ \ar{تحت}\\ $(\Delta)$};
\node at (12.675,-1.35){$+$};
\node[align=center,font=\small] at (12.675,-3.2){$(C_f)$\\ \ar{فوق}\\ $(\Delta)$};
\node at (9.05,-1.35){$0$};
\draw[densely dotted] (9.05,-0.9)--(9.05,-1.13) (9.05,-1.57)--(9.05,-1.8);
\draw (7.55,-4.6) .. controls (7.625,-3.06) and (8.6,-2.136) .. (9.05,-1.8) .. controls (9.5,-2.136) and (10.475,-3.06) .. (10.55,-4.6);
\node[align=center,font=\footnotesize,inner sep=0pt] at (9.05,-3.65){\ar{$(C_f)$ يقطع}\\ $(\Delta)$\\ \ar{في النقطة}\\ $A(0;-2)$};
\end{tikzpicture}
\end{FIG}
$(C_f)$ و $(\Delta)$ يتقاطعان في النقطة $(0;-2)$.

\QQ{2}
\pa $f$ قابلة للاشتقاق على $\R$ و $f'(x)=2e^{2x}-e^{x}-1$. ولدينا $\left(e^{x}-1\right)\left(2e^{x}+1\right)=2e^{2x}+e^{x}-2e^{x}-1=2e^{2x}-e^{x}-1$، إذن \bx{f'(x)=\left(e^{x}-1\right)\left(2e^{x}+1\right)}.

\pb $2e^{x}+1>0$، فإشارة $f'(x)$ هي إشارة $e^{x}-1$: سالبة على $]{-}\infty;0[$ وموجبة على $]0;+\infty[$. إذن \bxt{$f$ متناقصة تماما على $]{-}\infty;0]$ ومتزايدة تماما على $[0;+\infty[$}، و $f(0)=1-1-0-2=-2$.
\begin{TAB}
\tkzTabInit[lgt=2,espcl=3.5]{$x$/1,$f'(x)$/1,$f(x)$/2}{$-\infty$,$0$,$+\infty$}
\tkzTabLine{,-,z,+,}
\tkzTabVar{+/$+\infty$,-/$-2$,+/$+\infty$}
\end{TAB}

\QQ{3} $(T)$ يوازي $(\Delta)$ معناه $f'(x_0)=-1$:
\[2e^{2x_0}-e^{x_0}-1=-1\iff e^{x_0}\left(2e^{x_0}-1\right)=0\iff e^{x_0}=\frac12\iff x_0=-\ln2\]
(لأن $e^{x_0}>0$). و $f(-\ln2)=\dfrac14-\dfrac12+\ln2-2=\ln2-\dfrac94$، إذن:
$(T):\ y=-\left(x+\ln2\right)+\ln2-\dfrac94$، ومنه: \bx{(T):\ y=-x-\dfrac94}.

\QQ{4} \textbf{على $]{-}\infty;0]$:} $f$ مستمرة ومتناقصة تماما، و $f(-2{,}2)\approx0{,}10>0$، $f(-2{,}1)\approx-0{,}007<0$، إذن حسب مبرهنة القيم المتوسطة توجد قيمة وحيدة $\alpha\in]-2{,}2;-2{,}1[$ حيث $f(\alpha)=0$.

\textbf{على $[0;+\infty[$:} $f$ مستمرة ومتزايدة تماما، و $f(0{,}8)\approx-0{,}07<0$، $f(0{,}9)\approx0{,}69>0$، إذن توجد قيمة وحيدة $\beta\in]0{,}8;0{,}9[$ حيث $f(\beta)=0$.

وبما أن $f$ رتيبة تماما على كل من المجالين، فإن للمعادلة $f(x)=0$ \bxt{حلّين فقط: $-2{,}2<\alpha<-2{,}1$ و $0{,}8<\beta<0{,}9$}.

\QQ{5} $f''(x)=4e^{2x}-e^{x}=e^{x}\left(4e^{x}-1\right)$، وإشارتها إشارة $4e^{x}-1$، و $4e^{x}-1>0\iff x>-\ln4$.
إذن $f''$ تنعدم عند $-\ln4$ مغيّرة إشارتها، فالنقطة ذات الفاصلة $-\ln4=-2\ln2$ نقطة انعطاف، و:
\[f(-2\ln2)=e^{-4\ln2}-e^{-2\ln2}+2\ln2-2=\frac1{16}-\frac14+2\ln2-2=2\ln2-\frac{35}{16}\]
\[\bx{I\left(-2\ln2\,;\,2\ln2-\tfrac{35}{16}\right)}\qquad I(-1{,}39;-0{,}80)\]

\QQ{6}
\pa الرسم (الوحدة مصغّرة):
\begin{GR}{-4}{1.6}{-3}{3}{1.3}
\Ticks{-3,-2,-1,1}{-2,-1,1,2}
\draw[cThm,dashed,thick,domain=-4:1.6] plot(\x,{-\x-2}) node[pos=0.04,right]{$(\Delta)$};
\draw[cEx,thick,domain=-4:1.6] plot(\x,{-\x-2.25}) ;
\node[cEx] at (0.9,-2.7){$(T)$};
\draw[cDef,very thick,domain=-4:1.3,samples=150] plot(\x,{exp(2*\x)-exp(\x)-\x-2});
\fill(-2.107,0) circle(1.5pt) node[above]{$\alpha$};
\fill(0.811,0) circle(1.5pt) node[above left]{$\beta$};
\fill(-0.693,-1.557) circle(1.5pt);
\fill[cEx](-1.386,-0.801) circle(1.8pt) node[above right]{$I$};
\node[cDef] at (1.1,2.5){$(C_f)$};
\end{GR}

\pb $e^{2x}-e^{x}-m-2=0\iff e^{2x}-e^{x}-x-2=-x+m\iff f(x)=-x+m$.
الحلول هي فواصل نقط تقاطع $(C_f)$ مع المستقيم $(D_m):y=-x+m$ الموازي لـ $(\Delta)$ (ترتيبه إلى المبدأ $m$). المماس $(T)$ يوافق $m=-\dfrac94$ و $(\Delta)$ يوافق $m=-2$، و $(C_f)$ فوق $(T)$:
\begin{itemize}
\item $m<-\dfrac94$: لا توجد حلول.
\item $m=-\dfrac94$: حلّ وحيد مضاعف $x=-\ln2$.
\item $-\dfrac94<m<-2$: حلّان مختلفان.
\item $m\geq-2$: حلّ وحيد.
\end{itemize}
\end{hal}

% ===================== بكالوريا 2024 -- الموضوع الأول =====================
\begin{LTR}
\begin{tcolorbox}[enhanced,colback=cSec!8,colframe=cSec,arc=3mm,halign=center,boxrule=1pt,top=10pt,bottom=10pt,before upper={\selectlanguage{arabic}}]
{\color{cSec}\Large\bfseries هذه عيّنة مجانية من ملف «تمارين البكالوريا في الدالة الأسية»}\\[6pt]
الملف الكامل: كل مواضيع البكالوريا من 2008 إلى 2026 بالحلول المفصلة ، لكل الشعب العلمية\\[4pt]
وأكثر من 140 ملفًا آخر: ملخصات ، مذكرات ، تمارين نموذجية وتمارين البكالوريا بالحلول\\[4pt]
\foreignlanguage{english}{PDF -- Word -- LaTeX} قابلة للتعديل\\[8pt]
{\color{cThm}\Large\bfseries واتساب: \foreignlanguage{english}{0699 92 80 24}}
\end{tcolorbox}
\end{LTR}
\end{document}